\documentclass[11pt,a4paper]{article}

% ----- Packages for Styling and Formatting -----
\usepackage[utf8]{inputenc}
\usepackage[T1]{fontenc}
\usepackage{mathpazo} % Beautiful Palatino font
\usepackage[margin=0.9in]{geometry}
\usepackage{graphicx}
\usepackage{setspace}
\usepackage{amsmath}
\usepackage{booktabs}
\usepackage{xcolor}
\usepackage{fancyhdr}
\usepackage[colorlinks=true, linkcolor=blue, urlcolor=blue, citecolor=blue]{hyperref}
\usepackage{caption}
\usepackage{float}
\usepackage{titlesec}
\usepackage{minted} % For the python code in the appendix

% Define the custom blue color
\definecolor{CTtitle}{RGB}{30,90,160}

% Page styling for header
\pagestyle{fancy}
\fancyhf{}
\lhead{\textit{\textbf{CSG Final Project} --- April 2026}}
\rhead{\textit{\textbf{LLMs \& Dependency Grammar}}}
\renewcommand{\headrulewidth}{1pt}

% Formatting chapter headings nicely with blue rules
\titleformat{\section}
  {\color{CTtitle}\large\bfseries}{\thesection.}{0.5em}{}
  [\vspace{0.1cm}{\color{CTtitle}\titlerule[0.8pt]}]
\titleformat{\subsection}
  {\color{CTtitle}\normalsize\bfseries}{\thesubsection.}{0.5em}{}

\begin{document}

\begin{center}
    \vspace*{0.2cm}
    \textcolor{CTtitle}{\Large\textbf{Do Large Language Models Develop\\ Dependency Grammar?}}\\
    \vspace{0.2cm}
    \textit{\large Evidence from Attention-Derived Prefix Trees}
\end{center}
\vspace{0.4cm}
\begin{center}
    Jani Ravi Kailash (240486) \quad Aditya Panwar (240063) \quad Birkurwar Hitesh (240277)\\[0.15cm]
    Ishan Trikha (240471) \quad Sandeep Kumar Gupta (240928)\\[0.25cm]
    \small Computational Linguistics (CSG) --- Final Project \quad | \quad April 2026
\end{center}
\vspace{0.3cm}
{\color{CTtitle}\hrule height 1pt}
\vspace{0.3cm}
\begin{center}
    \textcolor{CTtitle}{\textbf{Abstract}}
\end{center}
\vspace{-0.25cm}
{\color{CTtitle}
We investigate whether transformer-based large language models (LLMs) internally develop dependency grammar without explicit supervision. Using attention matrices from \texttt{bert-base-multilingual-cased}, we derive incremental dependency trees over sentence prefixes and measure structural stability via \textit{Incremental Edge Change} (IEC). Across four typologically diverse languages---English, Hindi, German, and French---we find that LLM-derived trees are significantly more stable than random baselines ($p < 0.001$), that IEC decreases as sentences lengthen, and that middle BERT layers (5--8) exhibit the highest structural coherence. These findings suggest that transformers implicitly learn dependency-like representations through next-token prediction alone.
}
\vspace{0.3cm}
{\color{CTtitle}\hrule height 1pt}
\vspace{0.5cm}

\section{Motivation for the research problem}

Large Language Models (LLMs) such as BERT and GPT-2 demonstrated remarkable performance across a myriad of natural language processing tasks. However, the precise mechanisms by which they internally represented hierarchical sentence structures remained a deeply contested theoretical question. While bidirectional masked language models have access to full sentence contexts during training, human sentence processing uniquely operates in a strictly incremental, linear, and computationally constrained manner. Dependency grammar---a structural framework where every word is linked to a governing syntactic "head"---serves as the primary mechanism linguists use to model these compositional hierarchies.

Previous probing studies indicated that transformers extracted rich static syntactic trees for fully formed sentences. However, these static post-hoc analyses failed to capture the dynamic reality of language processing. As humans process incoming words from left to right, they construct mental dependency structures incrementally, integrating new words into a continuously maintained structural hypothesis. This process requires cognitive load management; massive restructuring of a sentence tree at every new word would overwhelm working memory. Thus, humans deploy strategies to stabilize the dependency grammar of the prefix string to avoid cognitive overload.

This project investigated whether transformer models inherently developed similar structural stability despite never receiving explicit topological supervision. We questioned whether the attention distributions in LLMs stabilized incrementally over time or whether they chaotically overhauled their entire dependency representations upon the arrival of every new token.

\textbf{Research objectives:} We possessed two core empirical goals:
\begin{enumerate}
    \item To investigate and mathematically quantify whether the dependency trees derived strictly from transformer attention mechanisms exhibited non-volatile, bounded structural updates as sentences grew incrementally longer.
    \item To strictly validate these structural behaviors across typologically diverse languages (English, Hindi, German, French) against theoretically derived random baseline graphs to confirm that emerging stability was a learned linguistic representation and not an arbitrary artifact of the transformer architecture.
\end{enumerate}

\section{Hypotheses and predictions}

We formulated two precisely testable hypotheses regarding the incremental structural topology of LLM attention distributions:

\textbf{Hypothesis 1 (Structural Stability):} The dependency arborescences extracted incrementally from the attention heads of the transformer would exhibit bounded mathematical changes. We measured this temporal volatility using Incremental Edge Change (IEC). We hypothesized that the IEC value would strictly decrease as the prefix length increased, demonstrating that the structural foundation of the sentence was stabilizing dynamically like a human cognitive structural parse.

\textbf{Hypothesis 2 (Statistical Non-Randomness):} If transformers were merely building graphs algorithmically without linguistic grounding, the topology changes would resemble random graph attachment. We hypothesized that the IEC derived from the transformer's attention matrices would be significantly and statistically lower than the IEC generated by memoryless, strictly random recursive tree structures at every corresponding prefix length.

\textbf{Predictions on the Data:}
\begin{itemize}
    \item The empirical IEC curves for the LLM would present a sharp downward trajectory as prefix length grew (bounded by a maximum prefix limit of 20 words), while the random baseline curve would remain highly volatile and roughly invariant over time.
    \item The intermediate layers of the multilingual BERT architecture (layers 5 through 8) would showcase the lowest volatility levels and the highest structural stability, adhering to prevailing linguistic probing literature identifying middle-stride layers as syntactic encoders.
    \item This non-volatile structural pattern is a generalized principle of predictive coding and would thus hold consistently across English, French, rigid German, and highly flexible Hindi dependency distributions.
\end{itemize}

\section{Methods}

We executed a comprehensive three-stage computational pipeline consisting of data preparation, attention extraction and aggregation, and graph evaluation.

\subsection{Data used}
We utilized corpora from the Universal Dependencies (UD v2.14) framework. To meet stringent sample size and diversity requirements, we sampled from four morphosyntactically distinct treebanks: English (EWT), Hindi (HDTB), German (GSD), and French (GSD). We programmatically filtered strictly for sentences containing between 5 and 20 tokens to observe sustained, incremental changes without interference from excessive truncation. We sampled exactly 100 sentences per language. These CoNLL-U format files provided the gold-standard (human-annotated) linguistic dependency trees serving as a benchmark for unlabeled attachment analysis. The data parsed into multi-layered python dictionaries tracking exact human-designed heads for each token.

\subsection{Extracting model attention}
We utilized the pretrained \texttt{bert-base-multilingual-cased} architecture to ensure identical parameter sets were responsible for parsing all four test languages. For a given sentence of length $n$, we fed the model strictly expanding prefixes $S_t$ such that $S_t = \{w_1, w_2, \dots, w_t\}$ for $t \in [2, n]$. 

At every incremental timestep $t$, we extracted the multi-headed attention matrices spanning all 12 operational layers. Because transformer architectures employ subword tokenization (WordPiece), we implemented a mathematically robust mapping layer: we aligned subword attention back to standard lexical boundaries by aggregating subword-to-subword correlations. For a word spanning subword indices $S$ and an attending head word spanning subword indices $T$, the aggregated scalar attention $A_{out}(u,v)$ was computed as the arithmetic mean across all mapped coordinates. We intentionally preserved the zero-indexed \texttt{[CLS]} token within the matrix; it functions as a global sink and an effective proxy for the linguistic "ROOT" of the dependency tree.

\subsection{Tree construction}
The inferred structural relationship relied on the premise that raw attention weights constituted directed edges in a syntactic graph. Specifically, elevated attention originating from word $A$ targeting word $B$ indicated a strong likelihood that $B$ served as the syntactic head governing $A$.

To cleanly translate dense $t \times t$ attention matrices into geometrically valid dependency trees without cycles, we deployed the Chu-Liu/Edmonds algorithm. This optimization located the maximum weight spanning arborescence rooted reliably at the designated \texttt{[CLS]} token node. The output of this pipeline stage was an explicit programmatic mapping delineating a single, valid governing head index for every active token present in the operational prefix.

\subsection{Inference method: testing structural volatility}
To formally validate our hypotheses, we invoked our primary inference methodology: measuring the \textbf{Incremental Edge Change (IEC)}. When the transition from $S_{t-1}$ to $S_t$ occurred via the introduction of token $w_t$, we compared the resulting arborescence $T_t$ against the predecessor $T_{t-1}$. We mathematically evaluated IEC as the fraction of previously observed tokens that were forced into a structural reassignment:

\begin{equation}
 IEC(t) = \frac{|\{ i \in \{1,\dots,t-1\} : head_t(i) \neq head_{t-1}(i) \}|}{t-1} 
\end{equation}

Under this inference framework, an IEC value of 0.0 reflects a perfectly stable syntactic expansion with zero structural volatility, while an IEC of 1.0 implies total structural collapse and hierarchical reorganization. We further evaluated absolute tree depth changes bounds to capture vertical volatility.

\textbf{Baseline Simulation and Statistical Test:}
To verify that structural stability was a learned language property, we synthesized recursive Monte-Carlo baseline graphs. We defined a random uniform attachment simulation where node $i$ strictly and uniformly sampled its parent from the set $\{0, 1, \dots, i-1\}$ where index 0 represented the ROOT. At every prefix length, we generated 200 random recursive paths and extracted expected stability bounds. To mathematically reject the null hypothesis across languages, we deployed a rigorous \textbf{Paired t-test} operating over the sequence of sample lengths, coupled with \textbf{Cohen's $d$} calculations to strictly evaluate the standardized effect size of structural coherence.

\section{Results}

Our analytical pipeline yielded high-density empirical measurements strictly favoring the hypothesized incremental learning models. Because initial processing identified layers 5 through 8 as structurally dominant, standard metric reporting concentrated distinctly on this intermediate grouping.

\subsection{Stability compared to random baseline}

\begin{figure}[h!]
    \centering
    \includegraphics[width=0.7\textwidth]{fig1_stability_curves.png}
    \caption{Structural Stability Curves: LLM vs Random Baseline.}
    \label{fig:stability}
\end{figure}

Consistent with Hypothesis 1, empirical testing demonstrated that attention-derived dependency trees grew significantly more structurally stagnant as sentence prefixes deepened. In contrast to high-volatility environments that frequently shuffle previous structural bounds, LLMs tightly bounded restructurings across all inputs.

Validating Hypothesis 2, the LLM stability vastly outperformed the simulated memoryless null-hypothesis graph networks. As calculated across the multi-lingual subsets, mean IEC limits for the transformer models reliably stabilized between 0.18 and 0.23 constraint bands. The paired t-test confirmed the statistical rejection of the uniform random attachment null hypothesis for every language cohort with strict confidence ($p < 0.001$). The observed Cohen's $d$ magnitude consistently ranged between 2.5 and 3.4, categorizing the structural coherence of transformer derivations as a massive, systematic effect far exceeding arbitrary bounds.

\subsection{Cross-lingual variations}

\begin{figure}[h!]
    \centering
    \includegraphics[width=0.7\textwidth]{fig2_language_comparison.png}
    \caption{Cross-Lingual Comparison: Mean IEC by Language.}
    \label{fig:crosslingual}
\end{figure}

Extending the inference evaluation across four structurally disparate linguistic families verified that incremental processing limits are foundational to the model, not merely superficial overfitting to rigid English phrase-structure architectures.

\begin{table}[h!]
    \centering
    \begin{tabular}{lcc}
    \toprule
    \textbf{Treebank Family} & \textbf{Mean IEC Limits} & \textbf{Random Baseline IEC} \\
    \midrule
    Hindi (HDTB)    & 0.187       & 0.553          \\
    English (EWT)   & 0.197       & 0.553          \\
    French (GSD)    & 0.215       & 0.553          \\
    German (GSD)    & 0.230       & 0.553          \\
    \bottomrule
    \end{tabular}
    \caption{Average IEC boundary analysis.}
    \label{tab:iec_scores}
\end{table}

Interestingly, the strictly head-final, morphologically rich language of Hindi effectively achieved the highest architectural stability score of 0.187. German displayed the highest volatility (0.230), probabilistically aligning with its complex separation of verbal components and inherently flexible word order constraints. However, no dataset crossed the 50\% threshold of the random volatility constraints.

\subsection{Layer-wise syntactic resolution}

\begin{figure}[h!]
    \centering
    \includegraphics[width=0.7\textwidth]{fig3_layer_analysis.png}
    \caption{Layer-wise Analysis of IEC.}
    \label{fig:layers}
\end{figure}

Extracting stability boundaries mapped against internal neural depth generated a distinct 'U-shaped' coherence topography. 
\begin{itemize}
    \item Early layers (1 through 3) operated erratically (IEC approximating 0.33 bounds), processing shallow n-gram lexical overlap constraints unconcerned with deep hierarchical relationships.
    \item Syntactic mid-layers (5 through 8) formed the stable computational core (IEC bounded between 0.15 and 0.19). We isolated deep grammatical parsing algorithms distinctly within this span. Unlabeled Attachment Scores (UAS) evaluated directly against human gold-standard trees further peaked distinctly within these intermediate modules.
    \item Final task layers (9 through 12) returned to volatile restructuring operations, aligning with known tendencies for upper layers to abandon rigid syntactic graphs to process holistic semantic pooling objectives.
\end{itemize}

\subsection{Internal Depth Volatility}

\begin{figure}[h!]
    \centering
    \includegraphics[width=0.7\textwidth]{fig4_depth_change.png}
    \caption{Tree Depth Volatility.}
    \label{fig:depth}
\end{figure}

Depth variation metrics corroborated edge changes. When a standard recursive network injected a novel item, maximum dependency depth predictably fluctuated erratically up to 2.5 distinct hierarchical leaps. In juxtaposition, LLMs bounded structural depth augmentations exclusively to minimal jumps of 0.5 to 1.0 elements. Such constrained topologies accurately simulate natural human dialect patterns, known empirically to avoid deep right-branching hierarchies through strategic node-skipping behaviors.

\section{Theoretical implications}

The documented implementation of incrementally stable graphs originating from massive neural networks carries profound theoretical implications spanning theories of machine intelligence algorithms to fundamental frameworks within evolutionary psycholinguistics. 

\textbf{Machine Structural Induction:}\\
Our investigation conclusively validates that explicit task scaffolding is profoundly unnecessary for structural grammatical grounding. Transformers, trained mathematically strictly on flat probabilistic linear next-token prediction, instinctively optimize their hidden geometrical topology to emulate deeply nested phrase-structure rules. The network dynamically learns that maintaining a non-volatile, relatively static "reference frame" of past head-relationships dramatically simplifies the computational gradient required for integrating novel word sequences.

\textbf{Implications for Cognitive Processing Theory:}\\
From a psycholinguistic perspective, these geometric structural consistencies precisely mirror the behavioral constraints of human cognitive processing. Gibson's \textbf{Dependency Locality Theory} explicitly argues that long-distance syntactic dependencies exponentially tax human working memory, causing integration costs to skyrocket. Humans combat this architectural bottleneck by rigidly committing to early branch processing, attempting to resolve incoming arguments against immediately accessible preceding head boundaries. 
Our metric inference demonstrates that transformers structurally behave identically to human memory load-reduction theories. Rather than recursively ripping down computational structures and rebuilding them per word---a technically permissible maneuver given bidirectional self-attention architectures---transformers organically learned that minimizing structural distance (volatility constraint) optimized prediction bounds.

\textbf{Language Evolution Dynamics:}\\
Finally, analyzing the stability metric aligns with the core tenet of the \textbf{Dependency Length Minimization (DLM)} principle of cross-lingual evolution in natural languages. Linguists establish that human lexical communication patterns universally evolve over time specifically to close the gap between disparate structural connections. Finding identical minimization paradigms naturally expressed across an AI---without prior evolutionary human bias introduced to the model---strongly concludes that non-volatile, shallow incremental tree structures are a deeply foundational consequence of optimal information-theoretic data traversal pathways.

% --- The Page Break directly enforces that references start on a new page, keeping text at 6 pages ---
\clearpage

\section*{References}
\begin{itemize}
    \item Chu, Y., \& Liu, T. (1965). On the shortest arborescence of a directed graph. \textit{Scientia Sinica}, 14, 1396–1400.
    \item Clark, K., Khandelwal, U., Levy, O., \& Manning, C. D. (2019). What does BERT look at? An analysis of BERT's attention. \textit{Proceedings of the 2019 ACL Workshop BlackboxNLP}.
    \item Devlin, J., Chang, M.-W., Lee, K., \& Toutanova, K. (2019). BERT: Pre-training of deep bidirectional transformers for language understanding. \textit{NAACL-HLT}.
    \item Edmonds, J. (1967). Optimum branchings. \textit{Journal of Research of the National Bureau of Standards}, 71B, 233–240.
    \item Futrell, R., Mahowald, K., \& Gibson, E. (2015). Large-scale evidence of dependency length minimization in 37 languages. \textit{PNAS}, 112(33), 10336–10341.
    \item Futrell, R., \& Levy, R. (2017). Noisy-context surprisal as a human sentence processing cost model. \textit{EACL}.
    \item Gibson, E. (2000). The dependency locality theory: A distance-based theory of linguistic complexity. In A. Marantz et al. (Eds.), \textit{Image, Language, Brain}, MIT Press.
    \item Hale, J. (2001). A probabilistic Earley parser as a psycholinguistic model. \textit{NAACL}.
    \item Hewitt, J., \& Manning, C. D. (2019). A structural probe for finding syntax in word representations. \textit{NAACL-HLT}.
    \item Jawahar, G., Sagot, B., \& Seddah, D. (2019). What does BERT learn about the structure of language? \textit{ACL}.
    \item Lewis, R. L., \& Vasishth, S. (2005). An activation-based model of sentence processing as skilled memory retrieval. \textit{Cognitive Science}, 29(3), 375–419.
    \item Mel'čuk, I. (1988). \textit{Dependency Syntax: Theory and Practice}. SUNY Press.
    \item Nivre, J. (2004). Incrementality in deterministic dependency parsing. \textit{Workshop on Incremental Parsing}.
    \item Nivre, J., de Marneffe, M.-C., Ginter, F., et al. (2020). Universal Dependencies v2: An evergrowing multilingual treebank collection. \textit{LREC}.
    \item Radford, A., Wu, J., Child, R., Luan, D., Amodei, D., \& Sutskever, I. (2019). Language models are unsupervised multitask learners. \textit{OpenAI Blog}.
    \item Sturt, P., \& Lombardo, V. (2005). Processing coordinated structures: Incrementality and connectedness. \textit{Cognitive Science}, 29(2), 291–305.
    \item Tenney, I., Das, D., \& Pavlick, E. (2019). BERT rediscovers the classical NLP pipeline. \textit{ACL}.
    \item Tesnière, L. (1959). \textit{Éléments de syntaxe structurale}. Klincksieck.
    \item Kovaleva, O., Romanov, A., Rogers, A., \& Rumshisky, A. (2019). Revealing the dark secrets of BERT. \textit{Proceedings of the 2019 Conference on Empirical Methods in Natural Language Processing (EMNLP)}.
    \item Voita, E., Talbot, D., Moiseev, F., Sennrich, R., \& Titov, I. (2019). Analyzing multi-head self-attention: Specialized heads do the heavy lifting, the rest can be pruned. \textit{ACL}.
\end{itemize}

\clearpage

\section*{Appendix A: Complete Python Implementation}

The complete, runnable pipeline is provided below. It is also publicly available on GitHub at:
\url{https://github.com/Sandeepgupta-24/LM1-dependency-grammar}

\vspace{0.5cm}

% Define a styled background for the syntax highlighted code
\definecolor{mintedbg}{HTML}{F8F9FA}

% Dynamically pull in the pipeline script, beautifully formatted
\inputminted[
    linenos=false,
    bgcolor=mintedbg,
    frame=leftline,
    framesep=3mm,
    framerule=1pt,
    rulecolor=\color{lightgray},
    fontsize=\footnotesize,
    tabsize=4,
    breaklines=true
]{python}{pipeline.py}

\end{document}
